\documentclass[12pt,a4paper]{article}

\usepackage[a4paper,margin=2cm]{geometry}
\usepackage[T2A]{fontenc}
\usepackage[utf8]{inputenc}
\usepackage[russian]{babel}
\usepackage{amsmath,amssymb,mathtools}
\usepackage{array,booktabs,longtable}
\usepackage{xcolor}
\usepackage{hyperref}
\usepackage{tikz}

\hypersetup{
  colorlinks=true,
  linkcolor=blue!55!black,
  urlcolor=blue!55!black,
  pdftitle={Ревью домашней работы},
  pdfauthor={Лёвин Артём Александрович}
}

\setlength{\parindent}{0pt}
\setlength{\parskip}{0.65em}
\renewcommand{\arraystretch}{1.25}

\begin{document}

\begin{titlepage}
\centering
\vspace*{0.22\textheight}
{\LARGE\bfseries Ревью домашней работы\par}
\vspace{2.2cm}
{\large Автор: Лёвин Артём Александрович\par}
\vspace{0.8cm}
{\large 16 сентября 2026 года\par}
\vfill
\begin{tikzpicture}
  \draw[line width=0.8pt] (0,0) .. controls (1.4,0.45) and (2.8,-0.45) .. (4.2,0)
                         .. controls (5.6,0.45) and (7.0,-0.45) .. (8.4,0);
\end{tikzpicture}
\end{titlepage}

\newpage
\section*{Сводная таблица}
\small
\setlength{\tabcolsep}{3pt}
\begin{longtable}{@{}>{\centering\arraybackslash}p{0.07\linewidth}p{0.17\linewidth}p{0.29\linewidth}p{0.20\linewidth}p{0.20\linewidth}@{}}
\toprule
№ задания & Тема & Суть ошибки & Ваш ответ & Правильный ответ \\
\midrule
\endfirsthead
\toprule
№ задания & Тема & Суть ошибки & Ваш ответ & Правильный ответ \\
\midrule
\endhead
\hyperref[task:1]{1} & Действия с десятичными и обыкновенными дробями & При сокращении дроби числитель и знаменатель изменены не одним и тем же множителем. & $6{,}27;\ 3{,}44;\ \frac{5}{12};\ \frac{2}{3}$ & $6{,}27;\ 3{,}44;\ \frac{5}{12};\ \frac{3}{4}$ \\
\addlinespace
\hyperref[task:2]{2} & Среднее арифметическое & После верной суммы первых четырёх чисел $26{,}0$ число $7$ прибавлено как десятая доля, поэтому получено $26{,}7$ вместо $33{,}0$. & В итоговой записи: $534$ & $6{,}6$ \\
\addlinespace
\hyperref[task:4]{4} & Среднее арифметическое двух чисел & Неверно составлена и преобразована зависимость между двумя числами; в числителе появляется $2y+22$, хотя разность должна оставаться равной $18$. & $31$ и $49$ & $75$ и $93$ \\
\addlinespace
\hyperref[task:5]{5} & Средняя скорость & Общий путь сложен неверно: $2{,}4+3{,}6+2$ записано как $6{,}2$ вместо $8$. Дополнительно неверно сокращено $\frac{30}{60}$. & $\frac{27960}{2000}$ км/ч & $16$ км/ч \\
\bottomrule
\end{longtable}
\normalsize

\newpage
\thispagestyle{empty}
\vspace*{\fill}
\begin{center}
{\LARGE\bfseries Разбор ошибок}
\end{center}
\vspace*{\fill}

\newpage
\phantomsection\label{task:1}
\section*{Задание 1}

\textbf{Текст задания.} Выполнить действия:
\[
3{,}47+2{,}8,\qquad 5{,}2-1{,}76,\qquad \frac{2}{3}-\frac{1}{4},\qquad \frac{18}{24}.
\]
Последнюю дробь сократить.

\textbf{Ваше решение.}
\[
3{,}47+2{,}8=6{,}27,
\qquad
5{,}20-1{,}76=3{,}44,
\]
\[
\frac{2}{3}-\frac{1}{4}
=\frac{8}{12}-\frac{3}{12}
=\frac{5}{12},
\qquad
\frac{18}{24}=\frac{2}{3}.
\]

\textbf{Ваш ответ.} $6{,}27;\ 3{,}44;\ \frac{5}{12};\ \frac{2}{3}$.

\textcolor{red}{\textbf{Ошибка:}} неверное сокращение дроби $\frac{18}{24}$.

\textbf{Где именно возникает ошибка.}
Последний корректный шаг в ошибочной части --- запись исходной дроби $\frac{18}{24}$. Первый неверный переход:
\[
\frac{18}{24}=\frac{2}{3}.
\]
Первые три действия в задании выполнены верно; ошибка появляется только при сокращении последней дроби.

\textbf{Почему этот шаг неверен.}
Сократить дробь означает разделить её числитель и знаменатель на одно и то же ненулевое число. Только при таком преобразовании значение дроби сохраняется. Для чисел $18$ и $24$ общий делитель равен $6$:
\[
18:6=3,\qquad 24:6=4.
\]
Поэтому правильное сокращение даёт $\frac{3}{4}$.

Получить из $18$ число $2$ можно делением на $9$, а из $24$ число $3$ --- делением на $8$. Делители разные, поэтому такое преобразование меняет значение дроби и не является сокращением. Это можно быстро проверить перекрёстным умножением: если бы дроби $\frac{18}{24}$ и $\frac{2}{3}$ были равны, произведения $18\cdot 3$ и $24\cdot 2$ совпали бы. Но
\[
18\cdot 3=54,\qquad 24\cdot 2=48,
\]
следовательно, дроби не равны.

\textcolor{blue!60!black}{\textbf{Ключевая идея:}} при сокращении дроби числитель и знаменатель обязательно делятся на один и тот же общий делитель.

\textbf{Исправление от последнего правильного шага.}
\[
\frac{18}{24}
=\frac{18:6}{24:6}
=\frac{3}{4}.
\]
Числа $3$ и $4$ уже не имеют общего делителя больше единицы, поэтому дробь сокращена полностью.

\textbf{Полное правильное решение.}
Первое действие:
\[
3{,}47+2{,}8=3{,}47+2{,}80=6{,}27.
\]
Второе действие:
\[
5{,}2-1{,}76=5{,}20-1{,}76=3{,}44.
\]
Для третьего действия приводим дроби к общему знаменателю $12$:
\[
\frac{2}{3}=\frac{8}{12},\qquad
\frac{1}{4}=\frac{3}{12},
\]
поэтому
\[
\frac{2}{3}-\frac{1}{4}
=\frac{8}{12}-\frac{3}{12}
=\frac{5}{12}.
\]
Последнюю дробь сокращаем на общий делитель $6$:
\[
\frac{18}{24}=\frac{3}{4}.
\]

\textbf{Проверка.}
Проверим только исправленную часть перекрёстным умножением:
\[
18\cdot 4=72,\qquad 24\cdot 3=72.
\]
Произведения равны, значит $\frac{18}{24}=\frac{3}{4}$.

\textcolor{teal!60!black}{\textbf{Итог:}} правильные результаты задания:
\[
6{,}27;\qquad 3{,}44;\qquad \frac{5}{12};\qquad \frac{3}{4}.
\]

\textbf{Задача на закрепление.} Сократите дробь $\frac{21}{28}$ до несократимого вида.

\newpage
\phantomsection\label{task:2}
\section*{Задание 2}

\textbf{Текст задания.} Найти среднее арифметическое чисел
\[
6{,}8;\qquad 7{,}2;\qquad 5{,}9;\qquad 6{,}1;\qquad 7.
\]

\textbf{Ваше решение.}
В записи сначала удобно сгруппированы пары $6{,}8+7{,}2$ и $5{,}9+6{,}1$. Получены промежуточные суммы $14{,}0$ и $12{,}0$, затем сумма первых четырёх чисел $26{,}0$. После этого в записи появляется $26{,}7$, а далее выполняется деление на $5$; справа в итоговой записи указано $534$.

\textbf{Ваш ответ.} В итоговой записи указано $534$.

\textcolor{red}{\textbf{Ошибка:}} неверное прибавление числа $7$ к промежуточной сумме $26{,}0$.

\textbf{Где именно возникает ошибка.}
Последний корректный шаг:
\[
(6{,}8+7{,}2)+(5{,}9+6{,}1)=14{,}0+12{,}0=26{,}0.
\]
Первый неверный шаг --- переход к числу $26{,}7$ после добавления последнего числа $7$.

\textbf{Почему этот шаг неверен.}
Число $7$ --- это семь целых, то есть $7{,}0$, а не семь десятых. Поэтому при сложении десятичные запятые нужно расположить друг под другом:
\[
26{,}0+7{,}0=33{,}0.
\]
Запись $26{,}7$ получилась бы при прибавлении $0{,}7$, но в условии дано число $7$.

Среднее арифметическое пяти чисел находят так: сначала складывают все пять чисел, затем полученную сумму делят на количество чисел, то есть на $5$. Если сумма найдена неверно, дальнейшее деление уже не может привести к правильному среднему.

\textcolor{blue!60!black}{\textbf{Ключевая идея:}} при сложении десятичных чисел нужно сохранять разряды; число $7$ при таком сложении следует мыслить как $7{,}0$.

\textbf{Исправление от последнего правильного шага.}
Продолжаем с верной суммы первых четырёх чисел:
\[
26{,}0+7{,}0=33{,}0.
\]
Теперь делим сумму пяти чисел на $5$:
\[
33{,}0:5=6{,}6.
\]

\textbf{Полное правильное решение.}
Удобно использовать предложенную в условии группировку:
\[
6{,}8+7{,}2=14{,}0,
\]
\[
5{,}9+6{,}1=12{,}0.
\]
Тогда сумма всех пяти чисел равна
\[
14{,}0+12{,}0+7{,}0=33{,}0.
\]
Чисел всего пять, поэтому среднее арифметическое равно
\[
\frac{33{,}0}{5}=6{,}6.
\]

\textbf{Проверка.}
Наименьшее из данных чисел равно $5{,}9$, наибольшее --- $7{,}2$. Среднее арифметическое должно находиться между ними. Число $6{,}6$ действительно удовлетворяет условию
\[
5{,}9<6{,}6<7{,}2.
\]
Кроме того,
\[
6{,}6\cdot 5=33,
\]
что возвращает найденную сумму всех чисел.

\textcolor{teal!60!black}{\textbf{Итог:}} среднее арифметическое данных пяти чисел равно $6{,}6$.

\textbf{Задача на закрепление.} Найдите среднее арифметическое чисел $4{,}8;\ 5{,}2;\ 6{,}7;\ 5{,}3;\ 8$.

\newpage
\phantomsection\label{task:4}
\section*{Задание 4}

\textbf{Текст задания.} Среднее арифметическое двух чисел равно $84$. Одно число больше другого на $18$. Найти оба числа.

\textbf{Ваше решение.}
В записи используется выражение
\[
\frac{x+(y+18)}{2}=84,
\]
после чего появляется переход
\[
\frac{2y+22}{2}=84.
\]
Далее записано
\[
2y+22=84,\qquad 2y=62,\qquad y=31,
\]
затем найдено второе число:
\[
31+18=49.
\]
В конце выполнена проверка полученных чисел:
\[
\frac{31+49}{2}=40.
\]

\textbf{Ваш ответ.} По вычислениям получены числа $31$ и $49$.

\textcolor{red}{\textbf{Ошибка:}} неверно составлена и преобразована связь между двумя числами.

\textbf{Где именно возникает ошибка.}
В первой строке одновременно используются две переменные $x$ и $y$, но не записана связь между ними. Первый однозначно неверный алгебраический переход:
\[
\frac{x+(y+18)}{2}=84
\quad\longrightarrow\quad
\frac{2y+22}{2}=84.
\]
Даже если предположить, что $x=y$, числитель должен был бы стать $2y+18$, а не $2y+22$.

\textbf{Почему этот шаг неверен.}
Условие сообщает только две вещи: среднее двух чисел равно $84$, а разность между большим и меньшим числом равна $18$. Удобнее ввести одну переменную. Если меньшее число обозначить через $y$, то большее число обязано быть равно $y+18$. Тогда оба числа выражены через одну переменную и условие о разности автоматически выполнено.

Среднее арифметическое двух чисел равно их сумме, делённой на $2$. Поэтому правильное уравнение имеет вид
\[
\frac{y+(y+18)}{2}=84.
\]
При раскрытии скобок в числителе получаем
\[
y+y+18=2y+18.
\]
Число $22$ здесь появиться не может: в условии задана прибавка $18$, и никакого преобразования, меняющего её на $22$, нет.

Есть и важная проверка: найденные в Вашем решении числа $31$ и $49$ действительно отличаются на $18$, но их среднее равно $40$, а не $84$. Значит, одно из двух условий задачи не выполнено.

\textcolor{blue!60!black}{\textbf{Ключевая идея:}} выразить оба искомых числа через одну переменную: меньшее $y$, большее $y+18$.

\textbf{Исправление.}
Составляем правильное уравнение:
\[
\frac{y+(y+18)}{2}=84.
\]
Умножаем обе части уравнения на $2$:
\[
y+y+18=168.
\]
Приводим подобные слагаемые:
\[
2y+18=168.
\]
Вычитаем $18$ из обеих частей:
\[
2y=150.
\]
Делим обе части на $2$:
\[
y=75.
\]
Большое число:
\[
75+18=93.
\]

\textbf{Полное правильное решение.}
Пусть меньшее число равно $y$. Тогда большее равно $y+18$. По условию их среднее арифметическое равно $84$:
\[
\frac{y+(y+18)}{2}=84.
\]
Умножаем на $2$:
\[
2y+18=168.
\]
Переносим $18$ в правую часть:
\[
2y=150.
\]
Отсюда
\[
y=75.
\]
Второе число равно
\[
75+18=93.
\]

\textbf{Проверка.}
Проверим оба условия задачи. Разность:
\[
93-75=18.
\]
Среднее арифметическое:
\[
\frac{75+93}{2}=\frac{168}{2}=84.
\]
Оба условия выполнены.

\textcolor{teal!60!black}{\textbf{Итог:}} искомые числа --- $75$ и $93$.

\textbf{Задача на закрепление.} Среднее арифметическое двух чисел равно $57$, а одно число больше другого на $14$. Найдите оба числа.

\newpage
\phantomsection\label{task:5}
\section*{Задание 5}

\textbf{Текст задания.} Велосипедист проехал первый участок длиной $2{,}4$ км за $8$ минут, второй участок длиной $3{,}6$ км за $12$ минут, третий участок длиной $2$ км за $10$ минут. Найти среднюю скорость велосипедиста на всём пути и выразить ответ в километрах в час.

\textbf{Ваше решение.}
В начале записаны отношения
\[
\frac{2{,}4}{8},\qquad \frac{3{,}6}{12},\qquad \frac{2}{10}.
\]
Затем вычисляется общий путь и общее время:
\[
2{,}4+3{,}6+2=6{,}2,
\]
\[
8+12+10=30.
\]
Далее в записи используется преобразование времени
\[
\frac{30}{60}=\frac{20}{30},
\]
после чего средняя скорость вычисляется через деление $6{,}2$ на полученную дробь. В конце записано
\[
\frac{27960}{2000}\ \text{км/ч}.
\]

\textbf{Ваш ответ.} $\frac{27960}{2000}$ км/ч.

\textcolor{red}{\textbf{Ошибка:}} неверно найден общий путь; дополнительно неверно преобразованы $30$ минут в часы.

\textbf{Где именно возникает ошибка.}
Правильная идея появляется в момент, когда отдельно складываются длины участков и отдельно времена движения. Первый неверный шаг:
\[
2{,}4+3{,}6+2=6{,}2.
\]
Сумма $2{,}4+3{,}6$ равна $6$, а после прибавления ещё $2$ получается $8$, а не $6{,}2$.

Позже есть ещё одна независимая ошибка:
\[
\frac{30}{60}=\frac{20}{30}.
\]
Дробь $\frac{30}{60}$ равна $\frac{1}{2}$, а $\frac{20}{30}$ равна $\frac{2}{3}$. Эти значения различны.

\textbf{Почему это неверно.}
Средняя скорость на всём пути определяется не как сумма скоростей на отдельных участках, а как отношение всего пройденного пути ко всему затраченному времени:
\[
v_{\text{ср}}=\frac{S_{\text{общ}}}{t_{\text{общ}}}.
\]
Поэтому сначала нужно безошибочно найти общий путь:
\[
2{,}4+3{,}6+2=6+2=8\ \text{км}.
\]
Общее время действительно найдено верно:
\[
8+12+10=30\ \text{минут}.
\]
Но ответ требуется в километрах в час, поэтому минуты необходимо перевести в часы. В одном часе $60$ минут, значит
\[
30\ \text{минут}=\frac{30}{60}\ \text{часа}=\frac{1}{2}\ \text{часа}.
\]
Сокращать дробь нужно делением числителя и знаменателя на одно и то же число. Делим $30$ и $60$ на $30$ и получаем $\frac12$.

\textcolor{blue!60!black}{\textbf{Ключевая идея:}} средняя скорость на всём маршруте равна общему пути, делённому на общее время; единицы времени перед делением должны соответствовать требуемой единице скорости.

\textbf{Исправление от последнего правильного шага.}
Сохраняем верную идею складывать пути и времена, но исправляем вычисления:
\[
S_{\text{общ}}=2{,}4+3{,}6+2=8\ \text{км},
\]
\[
t_{\text{общ}}=8+12+10=30\ \text{минут}=\frac12\ \text{часа}.
\]
Тогда
\[
v_{\text{ср}}=8:\frac12=8\cdot 2=16\ \text{км/ч}.
\]

\textbf{Полное правильное решение.}
Находим длину всего пути:
\[
S_{\text{общ}}=2{,}4+3{,}6+2=8\ \text{км}.
\]
Находим всё время движения:
\[
t_{\text{общ}}=8+12+10=30\ \text{минут}.
\]
Переводим время в часы:
\[
30\ \text{минут}=\frac{30}{60}\ \text{часа}=\frac12\ \text{часа}.
\]
Средняя скорость:
\[
v_{\text{ср}}=\frac{8}{1/2}=16\ \text{км/ч}.
\]

\textbf{Проверка.}
Если двигаться со средней скоростью $16$ км/ч в течение половины часа, пройденный путь будет
\[
16\cdot\frac12=8\ \text{км},
\]
что совпадает с общей длиной трёх участков. Значит, вычисление согласуется с условием.

\textcolor{teal!60!black}{\textbf{Итог:}} средняя скорость велосипедиста на всём пути равна $16$ км/ч.

\textbf{Задача на закрепление.} Велосипедист проехал $3$ км за $10$ минут, затем $4{,}5$ км за $15$ минут и ещё $2{,}5$ км за $5$ минут. Найдите среднюю скорость на всём пути в километрах в час.

\vfill
\begin{center}
\small Лёвин Артём Александрович эксклюзивно для Ксении Васильченко
\end{center}

\end{document}
